
\subsection{An obstruction to a direct leaf-layer recursion}
The path theorem does not license independent leaf reservations throughout a
tree. The following obstruction concerns a specified selection procedure;
it is not a lower bound for arbitrary mechanisms on trees.

Fix $0<\alpha<1$ and $0\le\beta<1$. In each orientation of the scaled
bipartite double cover, process left vertices in a leaf-layer order of the
original tree: all original leaves precede newly exposed leaves, and each
successive layer is exhausted before the next. Among maximum flows consistent
with previous choices, clamp each current left use to $\beta c_v$ within its
feasible interval. Pay owners by their resource fractions of these left uses,
as in \eqref{graph:prices}. On P4, $\alpha=1/\sqrt2$ and $\beta=1/2$
recover \eqref{graph:selector}. The collapsed-flow funding argument still
applies on every tree, but incentives need not.

\begin{proposition}[Failure of fixed leaf reserves]\label{graph:recursion-failure}
For every $0<\alpha<1$ and $0\le\beta<1$, this leaf-layer selector fails
full OUT on P6, already for a singleton coalition consisting of one mixed owner.
The failure is independent of tie-breaking within the first leaf layer and
of every selection choice after the first leaf in each orientation is fixed.
\end{proposition}
\begin{proof}
Use the path $0-1-2-3-4-5$, choose
\[
 0<q<\min\{1-\beta,\alpha^{-2}-1,1\},
 \qquad c(q)=(1,\alpha^{-1},q,q,\alpha^{-1},1).
\]
One fixed owner holds all resources at $0,2,3,5$; two outsiders hold resources
$1,4$ respectively. The first owner alone reduces both middle reports from
$q$ to $q/2$, leaving its two endpoint reports unchanged. Its withheld bundle
has production $q/2$ on edge $2-3$.

Consider the orientation with left vertices $\{0,2,4\}$. The flows
\[
 f_{0_L1_R}=1-q,\quad f_{2_L1_R}=q,\quad
 f_{4_L3_R}=\alpha q,\quad f_{4_L5_R}=\alpha
\]
are feasible: in particular $\alpha(1+q)<\alpha^{-1}$.
They saturate every right capacity, so their value $1+\alpha(1+q)$ is
maximum. Every maximum flow saturates $1_R$, whose only incoming sources
are $0_L,2_L$. Since the latter has capacity $q$, its original leaf satisfies
$w_0\ge1-q$. The displayed flow attains this lower endpoint. As
$1-q>\beta c_0=\beta$, the first leaf clamp fixes $w_0=1-q$.
Saturating $1_R$ then forces source $2_L$ to spend its entire $q$ there;
all other right use must come from $4_L$. Thus every remaining left use
is forced. Reflection proves the same statement for the other orientation,
whose first original leaf is $5_L$. Hence
\[
 w(q)=(1-q,\alpha(1+q),q,q,\alpha(1+q),1-q).
\]
The same argument applies at $q/2$, since all strict inequalities persist.
The deviating owner's total payment is
\[
 (1-q)+q+q+(1-q)=2
\]
at both reports. Its payment loss is zero, contradicting the required
outside-production compensation $q/2>0$.
\end{proof}

For the theorem's constants, $q=1/4$ gives the concrete capacities
$(1,\sqrt2,1/4,1/4,\sqrt2,1)$ and a productive withdrawal of $1/8$
at each middle vertex. Retuning a fixed fraction below one or taking
$\alpha$ closer to one cannot repair this leaf-layer architecture. The
proposition does not cover reserve fraction one, report-dependent reserves,
other global selectors, or orders that process a newly exposed vertex before
an original leaf in its orientation. In particular it establishes no
impossibility for arbitrary full-OUT tree mechanisms. A viable recursion
must retain additional information about the simultaneous payment response
across the central edge.
